% ==================================================
% 4. MANAGEMENT PROCESS -- PA2 REVISED
% Requires \PAChange, L column type, and longtable in main.tex.
% ==================================================

\section{Management Process}


% ==================================================
% 4.1 PROJECT ESTIMATES
% ==================================================

\subsection{Project Estimates}

\PAChange{The following estimates are planning
assumptions, not actual work hours.
The team has three members and five two-week
sprints. At 16 hours per person per week,
the total capacity is 480 person-hours.
A 20\% reserve leaves 384 planned person-hours.}

\begin{longtable}{|
    L{0.47\textwidth}|
    L{0.13\textwidth}|
    L{0.28\textwidth}|
}
\hline
\textbf{Workstream} &
\textbf{Estimate} &
\textbf{Lead / Support}
\\ \hline

\endfirsthead

\hline
\textbf{Workstream} &
\textbf{Estimate} &
\textbf{Lead / Support}
\\ \hline
\endhead

\PAChange{Requirements, UI planning, and design}
& \PAChange{44 h}
& \PAChange{Long / Trang / Thành}
\\ \hline

\PAChange{Medication data, mapping, and source checks}
& \PAChange{60 h}
& \PAChange{Long}
\\ \hline

\PAChange{OCR, intent NLP, voice, and manual input}
& \PAChange{60 h}
& \PAChange{Trang / Long / Thành}
\\ \hline

\PAChange{Medication safety checks and warnings}
& \PAChange{48 h}
& \PAChange{Long}
\\ \hline

\PAChange{HCP review, chatbot, and health summaries}
& \PAChange{42 h}
& \PAChange{Trang / Long}
\\ \hline

\PAChange{Reminders and adherence tracking}
& \PAChange{34 h}
& \PAChange{Thành / Long}
\\ \hline

\PAChange{Dataset review, accounts, and support features}
& \PAChange{30 h}
& \PAChange{Long / Trang / Thành}
\\ \hline

\PAChange{Integration, testing, fixes, and final documents}
& \PAChange{66 h}
& \PAChange{Whole team}
\\ \hline

\textbf{\PAChange{Total planned}}
& \textbf{\PAChange{384 h}}
& \PAChange{96 h reserved}
\\ \hline

\end{longtable}

\PAChange{Estimates cover student effort,
not paid labor. No new hardware or paid
services are currently budgeted.
The planned additional cash cost is 0 VND,
excluding existing equipment and Internet costs.}

\PAChange{The team reviews remaining effort
weekly and re-estimates when major requirements,
data availability, or technical issues affect
the plan. Approved changes are recorded
in GitHub Projects.}


% ==================================================
% 4.2 PROJECT PLAN
% ==================================================
\subsection{Project Plan}


% ==================================================
% 4.2.1 PHASE AND ITERATION PLAN
% ==================================================

\subsubsection{Phase and Iteration Plan}

The project follows three RUP phases:
Inception, Elaboration, and Construction,
organized into five two-week sprints.

\begin{longtable}{|
    L{0.13\textwidth}|
    L{0.13\textwidth}|
    L{0.17\textwidth}|
    L{0.23\textwidth}|
    L{0.24\textwidth}|
}
\hline

\textbf{Phase} &
\textbf{Sprint} &
\textbf{Period} &
\textbf{Objective} &
\textbf{Milestone / Achievement Criteria}
\\ \hline

\endfirsthead

\hline
\textbf{Phase} &
\textbf{Sprint} &
\textbf{Period} &
\textbf{Objective} &
\textbf{Milestone / Achievement Criteria}
\\ \hline

\endhead

Inception
& Sprint 1
& 5--18 Oct 2026
& Define the project scope, requirements,
team roles, main workflows, and initial system design.
& PA1 package completed; Vision Document,
initial architecture, UI prototype,
and project tracker prepared for review.
\\ \hline

Elaboration
& Sprint 2
& 19 Oct--1 Nov 2026
& Review and improve the existing requirements
and architecture; finalize module interfaces
and begin core feature development.
& Requirements and architecture reviewed;
module interfaces, data formats,
and test plans agreed by the team.
\\ \hline

Elaboration
& Sprint 3
& 2--15 Nov 2026
& Develop, integrate, and test medication input,
identification, and safety assessment.
& Four input methods demonstrated;
medication identification, interaction checks,
duplicate checks, and incomplete-data handling
tested with agreed test cases.
\\ \hline

Construction
& Sprint 4
& 16--29 Nov 2026
& Integrate the core Web application,
safety results, and simulated HCP review workflow.
& Integrated Web workflow, safety warnings,
and simulated HCP review pass agreed
integration tests.
\\ \hline

Construction
& Sprint 5
& 30 Nov--13 Dec 2026
& Complete medication reminders,
safety gates, adherence tracking,
and system testing.
& Core features pass acceptance tests;
critical defects are resolved;
final prototype, documentation,
and demo are ready.
\\ \hline

\end{longtable}



% ==================================================
% 4.2.2 ITERATION OBJECTIVES
% ==================================================

\subsubsection{Iteration Objectives}

\begin{longtable}{|
    L{0.13\textwidth}|
    L{0.55\textwidth}|
    L{0.20\textwidth}|
}
\hline

\textbf{Sprint} &
\textbf{Iteration Objectives} &
\textbf{Owner(s)}
\\ \hline

\endfirsthead

\hline
\textbf{Sprint} &
\textbf{Iteration Objectives} &
\textbf{Owner(s)}
\\ \hline

\endhead

% --------------------------------------------------
% SPRINT 1
% --------------------------------------------------

Sprint 1
& Define the project scope and main requirements;
prepare the initial architecture, workflows,
UI prototype, Project Plan, and Vision Document.
Set up project tracking and team responsibilities.
& Whole team
\\ \hline

% --------------------------------------------------
% SPRINT 2
% --------------------------------------------------

Sprint 2
& Review and improve the existing requirements
and architecture. Finalize module interfaces,
data formats, safety rules, and test plans.
Begin developing the main system features.

\PAChange{Revise the Vision Document and Project Plan;
prepare 18 use cases, corresponding diagrams,
and UI flows. Review medication data readiness
and integration requirements.}
& Whole team
\\ \hline

% --------------------------------------------------
% SPRINT 3
% --------------------------------------------------

Sprint 3
& Long develops the NLP chatbot and improves
medicine data processing, identification,
and safety checks (Modules 2 and 3).

Trang develops OCR and user confirmation.
Thành develops Voice and Manual Input
with medication suggestions and auto-fill.
Integrate and test the core modules.

\PAChange{Test supported drug interactions,
duplicate ingredients, missing-data handling,
and severe interaction warnings using
controlled test cases.}
& Whole team
\\ \hline

% --------------------------------------------------
% SPRINT 4
% --------------------------------------------------

Sprint 4
& Build and integrate the main Web application.
Connect medication input, identification,
safety checks, and warnings.
Trang develops the simulated HCP review
workflow (Module 4).
Test the main end-to-end workflows.

\PAChange{Add patient--HCP case conversations,
medication history, chatbot support,
patient health summaries, and RBAC.
Begin dataset review, administration,
and technical support workflows.}
& Whole team
\\ \hline

% --------------------------------------------------
% SPRINT 5
% --------------------------------------------------

Sprint 5
& Thành completes medication reminders,
safety gates, and adherence tracking (Module 5).
The whole team performs regression testing,
fixes critical issues, and prepares the final
documentation and demonstration.

\PAChange{Complete the supporting workflows
started in Sprint 4. Test reminder restrictions,
access permissions, and all agreed
current-scope use cases before release.}
& Whole team
\\ \hline

\end{longtable}



% ==================================================
% 4.2.3 BUILDS AND RELEASES
% ==================================================

\subsubsection{Builds and Releases}

\begin{longtable}{|
    L{0.17\textwidth}|
    L{0.16\textwidth}|
    L{0.17\textwidth}|
    L{0.38\textwidth}|
}
\hline
\textbf{Build} &
\textbf{Target Date} &
\textbf{Type} &
\textbf{Acceptance Check}
\\ \hline

\endfirsthead

\hline
\textbf{Build} &
\textbf{Target Date} &
\textbf{Type} &
\textbf{Acceptance Check}
\\ \hline
\endhead

UI Prototype
& 18 Oct 2026
& Design prototype
& Initial UI screens and main user workflows
are prepared for review.
\\ \hline

\PAChange{PA2 Baseline}
& \PAChange{1 Nov 2026}
& \PAChange{Requirements}
& \PAChange{Revised Vision and Project Plan,
18 use-case specifications and diagrams,
and updated UI designs are prepared
for review.}
\\ \hline

Build 0.1
& 15 Nov 2026
& Technical PoC
& \PAChange{Core medication intake,
identification, and safety checks
are demonstrated and tested.
Unsupported DDI checks are reported
as Unassessed. Controlled test cases
are used when verified mappings
are unavailable.}
\\ \hline

Build 0.2
& 29 Nov 2026
& Integrated test
& \PAChange{The main Web workflow,
safety results, and simulated HCP
review pass integration tests.
Patient--HCP conversations,
health summaries, and access controls
are tested within supported scope.
Initial supporting workflows are available
for testing.}
\\ \hline

Build 1.0
& 13 Dec 2026
& Final prototype
& \PAChange{Medication reminders,
safety gates, adherence tracking,
and remaining current-scope workflows
pass agreed acceptance tests.
Critical defects are resolved,
and final documentation and
demo materials are prepared.}
\\ \hline

\end{longtable}



% ==================================================
% 4.2.4 PROJECT SCHEDULE
% ==================================================

\subsubsection{Project Schedule}

\PAChange{The following schedule keeps the original
PA1 milestones and includes the updated PA2 scope.
Actual progress and supporting evidence are
tracked through GitHub and weekly reports.}

\begin{longtable}{|
    L{0.12\textwidth}|
    L{0.15\textwidth}|
    L{0.35\textwidth}|
    L{0.14\textwidth}|
    L{0.14\textwidth}|
}
\hline

\textbf{Period} &
\textbf{Workstream} &
\textbf{Key Tasks} &
\textbf{Owner(s)} &
\textbf{Output}
\\ \hline

\endfirsthead

\hline
\textbf{Period} &
\textbf{Workstream} &
\textbf{Key Tasks} &
\textbf{Owner(s)} &
\textbf{Output}
\\ \hline

\endhead

% --------------------------------------------------
% SPRINT 1
% --------------------------------------------------

5--11 Oct
& Project definition
& Define the problem, scope, main features,
team roles, data sources, and initial
requirements. Prepare the PA0 proposal.
& Whole team
& PA0 and scope
\\ \hline

12--18 Oct
& PA1 planning
& Complete the Project Plan, Vision Document,
initial architecture, main workflows,
UI prototype, risk list, and weekly reports.
Set up project tracking.
& Whole team
& PA1 package
\\ \hline

% --------------------------------------------------
% SPRINT 2
% --------------------------------------------------

19--25 Oct
& Design review
& Review existing requirements, workflows,
and architecture. Finalize user confirmation,
data formats, module interfaces,
and safety rules.
& Whole team
& Reviewed design
\\ \hline

26 Oct--1 Nov
& Core preparation
& Long reviews medication data and
DDInter mappings. Trang improves OCR,
UI/UX, and HCP design. Thành designs
Voice, Manual Input, and reminder flows.
Prepare module tests and integration
interfaces.

\PAChange{Revise the Project Plan and Vision
Document; prepare 18 use cases and diagrams;
update UI flows and review data readiness.}
& Long /
Trang /
Thành
& \PAChange{PA2 requirements
and design baseline}
\\ \hline

% --------------------------------------------------
% SPRINT 3
% --------------------------------------------------

2--8 Nov
& Module development
& Long develops the NLP chatbot and
improves Modules 2 and 3.
Trang develops OCR and confirmation
screens. Thành develops Voice and
Manual Input with medication suggestions
and auto-fill.
& Long /
Trang /
Thành
& Working modules
\\ \hline

9--15 Nov
& PoC integration
& Connect the four input methods to
medication identification and safety
assessment. Test DDInter 2.0 mapping,
duplicate ingredients, uncertain results,
and user confirmation.

\PAChange{Test severe warnings and
missing-data handling. Use controlled
test cases when verified DDI mappings
are unavailable.}
& Whole team
& Build 0.1
\\ \hline

% --------------------------------------------------
% SPRINT 4
% --------------------------------------------------

16--22 Nov
& Core Web integration
& Build the main Web pages and APIs.
Integrate medication search, safety
checks, warnings, and result display.
Trang develops the simulated HCP
review workflow.

\PAChange{Add HCP case conversations,
chatbot support, health summaries,
and RBAC. Begin dataset review and
administrative workflows.}
& Whole team
& Integrated alpha
\\ \hline

23--29 Nov
& Integration testing
& Complete HCP review integration,
test end-to-end workflows,
handle missing data and errors,
and fix major integration issues.

\PAChange{Test case permissions,
health summary sharing,
medication history, and review outcomes.
Begin technical support workflows.}
& Whole team
& Build 0.2
\\ \hline

% --------------------------------------------------
% SPRINT 5
% --------------------------------------------------

30 Nov--6 Dec
& Reminders
& Thành completes medication reminders,
the safety gate, and adherence tracking.
Long and Trang support data integration,
interface testing, and safety review.

\PAChange{Complete dataset review,
account management, administrative
history, and technical support.
Test reminder restrictions and RBAC.}
& Thành /
Team
& Reminder beta
\\ \hline

7--13 Dec
& Final testing
& Run regression, usability, and
acceptance tests. Fix critical issues,
prepare documentation, release notes,
and the final demonstration.

\PAChange{Verify the agreed current-scope
use cases and safety requirements.}
& Whole team
& Build 1.0
\\ \hline

\end{longtable}





% ==================================================
% 4.2.5 PROJECT RESOURCING
% ==================================================

\subsubsection{Project Resourcing}

\begin{longtable}{|
    L{0.22\textwidth}|
    L{0.32\textwidth}|
    L{0.34\textwidth}|
}
\hline

\textbf{Phase / Sprint} &
\textbf{People and Responsibilities} &
\textbf{Skills and Resources}
\\ \hline

\endfirsthead

\hline
\textbf{Phase / Sprint} &
\textbf{People and Responsibilities} &
\textbf{Skills and Resources}
\\ \hline
\endhead

% SPRINT 1

Inception / Sprint 1
& Whole team: planning, requirements,
initial architecture, and UI prototype.
& Git, LaTeX, UI design, and
requirements analysis.
\\ \hline

% SPRINT 2

Elaboration / Sprint 2
& Long: NLP, medication data, and
safety design.

Trang: OCR, UI/UX, and HCP design.

Thành: Voice, Manual Input,
and reminder design.

\PAChange{Whole team: PA2 requirements,
use cases, and UI updates.}
& Data preprocessing, DDInter 2.0,
OCR, NLP, speech-to-text,
API design, and safety requirements.

\PAChange{Medication data validation
and use-case modeling.}
\\ \hline

% SPRINT 3

Elaboration / Sprint 3
& Long: NLP, Module 2,
and Module 3.

Trang: OCR and UI confirmation.

Thành: Voice and Manual Input.

Whole team: integration and testing.
& Python, medication datasets,
DDInter 2.0, OCR tools,
speech tools, API integration,
and automated tests.

\PAChange{Controlled DDI test cases
and missing-data handling.}
\\ \hline

% SPRINT 4

Construction / Sprint 4
& Trang: UI/UX and Module 4.

Long: safety service integration.

Thành: input integration.

Whole team: system testing.

\PAChange{Develop HCP conversations,
health summaries, RBAC, and initial
supporting workflows.}
& Web development, API integration,
HCP workflow, testing tools,
and issue tracking.

\PAChange{Authentication, access control,
and integration testing.}
\\ \hline

% SPRINT 5

Construction / Sprint 5
& Thành: Module 5.

Long and Trang: integration
and safety review.

Whole team: testing and final release.

\PAChange{Complete dataset review,
administration, and technical support
workflows.}
& Reminder logic, safety gate testing,
regression testing, documentation,
and presentation tools.

\PAChange{Role-based access testing
and final acceptance tests.}
\\ \hline

\end{longtable}

The project has a fixed team of three students,
so no additional recruitment is planned.
Team members will learn the required tools
and technologies based on their assigned tasks.

Training on OCR, NLP, speech-to-text,
medication data, and DDInter 2.0
will be completed as needed before Sprint 3.
Knowledge sharing and technical support
will continue throughout the project.

\PAChange{The team uses GitHub for source code,
collaboration, and project tracking.
Medication data sources include DAV,
RxNorm, WHO INN, DDInter 2.0,
DailyMed, openFDA, and other collected
medical and manufacturer references.
WIPO Lex supports data-rights review.}

\PAChange{Collected sources support research
and data preprocessing. Medication mappings
and clinical information require validation
before use in safety assessments.
Testing uses synthetic
or de-identified data.}




% ==================================================
% 4.3 PROJECT MONITORING AND CONTROL
% ==================================================
\subsection{Project Monitoring and Control}

% 4.3.1 -- Change control
\subsubsection{Requirements Management}

The Vision Document defines the main project requirements.
Detailed requirements are managed through use cases
and GitHub Issues. Each requirement has a priority,
assigned owner, and acceptance criteria.
Safety-related requirements must also cover possible
errors and missing data.

\PAChange{The PA2 baseline contains 13 product features,
10 non-functional requirements, and 18 use cases.
Changes are linked to affected use cases, UI screens,
tests, and project tasks.}

Requirement changes are recorded in GitHub Issues
and reviewed based on their impact on project scope,
safety, and schedule. The Project Manager and the
responsible team member review each change.
Major changes require agreement from the whole team.
Approved changes are updated in the related documents
and GitHub Projects.

Requirements and task progress are tracked using
GitHub Projects and reviewed during each sprint.
A requirement is considered complete when its
acceptance criteria are met and verified through
appropriate reviews or tests.

% 4.3.2 -- Reporting and numbers
\subsubsection{Reporting and Measurement}

\PAChange{GitHub is the team's central workspace for
project execution and tracking. GitHub Projects organizes
sprint tasks, owners, deadlines, status, and blockers;
GitHub Issues records tasks and defects. Members update
their assigned work each week, and the Project Manager
uses these records to prepare weekly status reports.}

\begin{longtable}{|L{0.34\textwidth}|L{0.54\textwidth}|}
\hline
\textbf{Measure} & \textbf{How It Is Checked} \\ \hline
\endfirsthead
\hline
\textbf{Measure} & \textbf{How It Is Checked} \\ \hline
\endhead
\PAChange{Task progress} & \PAChange{Completed tasks versus planned tasks for the sprint; incomplete tasks keep an owner and next action.} \\ \hline
\PAChange{Milestone status} & \PAChange{Compare each planned deliverable with its acceptance check and target date.} \\ \hline
\PAChange{Test results} & \PAChange{Passed, failed, and blocked test cases; track critical defects until retested.} \\ \hline
\PAChange{Data readiness} & \PAChange{Report accepted medicine records, usable DDI mappings, and unsupported checks separately.} \\ \hline
\PAChange{Risk and workload} & \PAChange{Review open high risks, blockers, and estimated work left each week.} \\ \hline
\PAChange{Team collaboration} & \PAChange{Review linked Issues, commits, and pull requests to check contributions and code reviews.} \\ \hline
\end{longtable}

% PA2 -- Evidence of execution and teamwork
\PAChange{Issues, commits, pull requests, and test results
provide evidence of completed work and team reviews.
Tasks are marked Completed only after their acceptance
criteria are checked. Weekly reports summarize progress,
contributions, delays, and next actions.}

\PAChange{A missed milestone or critical safety defect is
reported to the whole team. The Project Manager
reassigns work or reduces non-essential scope;
mandatory safety rules are not removed.}


% ==================================================
% 4.3.3 RISK MANAGEMENT
% ==================================================

\subsubsection{Risk Management}

Project risks are identified during planning,
development, and testing. Each risk is ranked
by likelihood and impact, with an assigned
owner and a mitigation plan.

The team reviews risks weekly and tracks
related issues using GitHub Issues.
Risks are updated based on project
progress and test results.

\begin{longtable}{|
    L{0.07\textwidth}|
    L{0.10\textwidth}|
    L{0.25\textwidth}|
    L{0.29\textwidth}|
    L{0.10\textwidth}|
    L{0.07\textwidth}|
}
\hline

\textbf{ID} &
\textbf{Ranking} &
\textbf{Description and Impact} &
\textbf{Mitigation / Contingency} &
\textbf{Owner} &
\textbf{Status}
\\ \hline

\endfirsthead

\hline
\textbf{ID} &
\textbf{Ranking} &
\textbf{Description and Impact} &
\textbf{Mitigation / Contingency} &
\textbf{Owner} &
\textbf{Status}
\\ \hline

\endhead

% R-01: MEDICATION IDENTIFICATION

R-01
& High
& Incorrect medication identification
or ingredient mapping may lead to
wrong safety results.
& Verify medication details, review mappings,
and keep unclear matches unresolved.
& Long
& Open
\\ \hline

% R-02: OCR ERRORS

R-02
& High
& OCR errors may cause incorrect
medication names, strengths,
or schedules.
& Require user confirmation and provide
manual correction or entry.
& Trang
& Open
\\ \hline

% R-03: CHATBOT SAFETY

R-03
& High
& The NLP chatbot may produce
incorrect or unsupported
medical explanations.
& Base explanations on verified
safety results, limit AI decisions,
and test unsafe responses.
& Long
& Open
\\ \hline

% R-04: VOICE RECOGNITION

R-04
& High
& Voice recognition may misunderstand
medication names or instructions.
& Show editable transcripts,
require confirmation, and provide
Manual Input as a fallback.
& Thành
& Open
\\ \hline

% R-05: DATA COVERAGE AND RIGHTS

R-05
& High
& Limited medication data,
DDInter coverage, or data licensing
may affect safety checks.
& Use permitted data, record source
limits, and support replaceable
data sources.

\PAChange{Keep unsupported DDI checks
Unassessed and use controlled test
cases when verified mappings
are unavailable.}
& Long
& Open
\\ \hline

% R-06: MISINTERPRETED SAFETY RESULTS

R-06
& High
& Users may believe that no detected
interaction means the medications
are completely safe.
& Show clear risk levels,
data limitations, and warnings
when information is incomplete.

\PAChange{Clearly distinguish assessed,
partially assessed, unassessed,
and failed results.}
& Trang /
Long
& Open
\\ \hline

% R-07: PRIVACY

R-07
& Medium
& Sensitive medication and personal
information may be exposed.
& Use test data, limit access,
protect user information, and avoid
storing sensitive data unnecessarily.

\PAChange{Apply RBAC and require patient
permission for HCP case access.}
& Whole team
& Open
\\ \hline

% R-08: REMINDER SAFETY GATE

R-08
& High
& The reminder safety gate may fail
to block new reminders for drugs

with \PAChange{severe interactions} or may block
unrelated medications.

& Test the gate for each affected
medication. Keep existing reminders
unchanged and require HCP review
before activating affected
new reminders.

\PAChange{Preserve detected severe
warnings if later checks fail.}
& Thành
& Open
\\ \hline

% R-09: TIME AND WORKLOAD

R-09
& Medium
& Limited time and uneven workload
may delay project completion.
& Prioritize core features,
track tasks weekly, and adjust
workloads when needed.
& Trang
& Open
\\ \hline

% R-10: MODULE INTEGRATION

R-10
& Medium
& Different module data formats
may cause integration errors.
& Define a common medication
data format, agree on module
interfaces, and test integration early.
& Whole team
& Open
\\ \hline

% R-11: SIMULATED HCP REVIEW

R-11
& Medium
& Simulated HCP review may not
reflect real medical practice,
and corrections may leave
outdated safety results.
& Clearly mark HCP review
as simulated. Record review
decisions and reassess affected
data after corrections.

\PAChange{Keep case conversations
and decisions linked to
their review records.}
& Trang
& Open
\\ \hline

% R-12: NOTIFICATION RELIABILITY

R-12
& Medium
& Reminder notifications may be
delayed, missed, or sent
more than once.
& Track delivery status,
test notification failures,
and prevent duplicate reminders.

\PAChange{Keep reminder schedules
available inside the application
when notifications fail.}
& Thành
& Open
\\ \hline

% R-13: UNAUTHORIZED ACCESS - PA2 ADDITION

\PAChange{R-13}
& \PAChange{High}
& \PAChange{Unauthorized users may
access HCP review cases or
restricted patient information.}
& \PAChange{Test role permissions
and case-level access.
Deny access when permission
cannot be verified.}
& \PAChange{Trang}
& \PAChange{Open}
\\ \hline

\end{longtable}


% 4.3.4 -- Version control, quality gate, and data release
\subsubsection{Configuration Management}

\PAChange{The team manages code, documents, and test assets
in a shared GitHub repository. Work is done on separate
branches and merged through reviewed pull requests.}

\PAChange{\textbf{Project repository:}
\href{https://github.com/hyt2434/Medical_Safety_Assistant}{\texttt{hyt2434/Medical\_Safety\_Assistant}}.
The repository is public; restricted datasets, secrets,
and identifiable patient information must not be published.}

\begin{longtable}{|L{0.31\textwidth}|L{0.57\textwidth}|}
\hline
\textbf{Item} & \textbf{Control Rule} \\ \hline
\endfirsthead
\hline
\textbf{Item} & \textbf{Control Rule} \\ \hline
\endhead
\PAChange{Source code} & \PAChange{Use feature branches and pull requests; check that changed features build and pass relevant tests before merging into the main branch.} \\ \hline
\PAChange{Requirements and design} & \PAChange{Version the Vision, Project Plan, use cases, and UI files; record major edits in Revision History.} \\ \hline
\PAChange{Test evidence} & \PAChange{Store test cases, results, and defect references with the related build or issue.} \\ \hline
\PAChange{Medication datasets} & \PAChange{Track source, version, audit status, and allowed use. Do not publish restricted or unreviewed data.} \\ \hline
\PAChange{Secrets and patient data} & \PAChange{Do not commit credentials, access tokens, or identifiable patient data; use synthetic or de-identified test records.} \\ \hline
\PAChange{Builds and rollback} & \PAChange{Tag reviewed builds 0.1, 0.2, and 1.0; keep earlier working versions for recovery if an update fails.} \\ \hline
\end{longtable}

\PAChange{Trang coordinates merges and release records.
All members review changes in their assigned areas.}
